% Self-contained manuscript. Compile twice with pdfLaTeX.
% All proofs and bibliography entries are included in this file.
\documentclass[11pt]{article}
\usepackage[T1]{fontenc}
\usepackage[utf8]{inputenc}
\usepackage{lmodern,amsmath,amssymb,amsthm,mathtools}
\usepackage[margin=1in]{geometry}
\usepackage{microtype}
\usepackage{enumitem,booktabs,tabularx,array}
\usepackage{caption}
\usepackage{xurl}
\usepackage[hidelinks]{hyperref}
\usepackage[nameinlink,noabbrev]{cleveref}
\newcommand{\R}{\mathbb R}
\newcommand{\E}{\mathbb E}
\newcommand{\Pp}{\mathbb P}
\newcommand{\norm}[1]{\left\lVert #1\right\rVert}
\newcommand{\ip}[2]{\left\langle #1,#2\right\rangle}
\newcommand{\eps}{\varepsilon}
\newcommand{\Lb}{\bar L}
\newcommand{\KL}{\operatorname{KL}}
\newcommand{\Tstar}{\mathsf T_1^\star}
\newtheorem{theorem}{Theorem}[section]
\newtheorem{lemma}[theorem]{Lemma}
\newtheorem{proposition}[theorem]{Proposition}
\newtheorem{corollary}[theorem]{Corollary}
\theoremstyle{definition}
\newtheorem{definition}[theorem]{Definition}
\newtheorem{assumption}[theorem]{Assumption}
\theoremstyle{remark}
\newtheorem{remark}[theorem]{Remark}
\setlength{\parindent}{15pt}
\setlength{\parskip}{3pt}
\setlength{\emergencystretch}{2em}
\setlist[itemize]{leftmargin=*,topsep=4pt,itemsep=3pt}
\setlist[enumerate]{leftmargin=*,topsep=4pt,itemsep=3pt}
\captionsetup{font=small,labelfont=bf}
\allowdisplaybreaks[2]
\hypersetup{
  pdftitle={Sharp Single-Query Complexity of Nonconvex Optimization under Mean-Squared Smoothness},
  pdfauthor={Colombo Team},
  pdfsubject={Minimax oracle complexity of stochastic nonconvex optimization},
  pdfkeywords={stochastic nonconvex optimization, mean-squared smoothness, oracle complexity, sample reuse}
}

\title{\textbf{Sharp Single-Query Complexity of Nonconvex Optimization under Mean-Squared Smoothness}}
\author{Colombo Team\\[0.4em]
\small\href{mailto:colombo.2026@outlook.com}{\texttt{colombo.2026@outlook.com}}}
\date{}

\begin{document}
\maketitle

\begin{abstract}
Mean-squared smoothness supports variance reduction when the same random
sample can be evaluated at two points. We show that, when each fresh sample
is evaluated at only one point, stochastic gradient descent is minimax
optimal despite this additional regularity. For initial objective gap
$\Delta$, mean-squared smoothness constant $\bar L$, and gradient-noise
variance bound $\sigma^2$, the dimension-uniform query complexity of finding
$\widehat x$ with $\E\|\nabla F(\widehat x)\|\le\eps$ is
\[
 \mathsf T_1^\star(\Delta,\bar L,\sigma,\eps)
 =\Theta\!\left(
   \frac{\Delta\bar L}{\eps^2}
   +\frac{\Delta\bar L\sigma^2}{\eps^4}
 \right),
 \qquad 0<\eps\le c_0\sqrt{\Delta\bar L},
\]
with universal constants. The lower bound holds for arbitrary adaptive
randomized queries and unqueried outputs on an unbounded domain. It is
realized by scalar sample losses whose gradients are globally
$\bar L$-Lipschitz. The construction combines a sequence of smooth gates
with additive Gaussian noise in a shared variable block. Completing each
stage requires learning a hidden direction. Exact gradients in earlier
blocks can reveal scalar projections of that direction, so the information
argument controls both these projections and the noisy vector response.
A sequential coupling then charges a separate estimation cost at every
stage. Together with the classical SGD upper bound, the result establishes
a separation between $\Theta(\eps^{-4})$ single-query complexity and
$\Theta(\eps^{-3})$ complexity with two evaluations per sample, for fixed
positive problem parameters.
\end{abstract}

\noindent\textbf{Keywords:} Stochastic nonconvex optimization; oracle
complexity; mean-squared smoothness; randomized lower bounds; sample reuse.
\par\medskip
\begin{samepage}
\noindent\textbf{AI Usage.}
Nearly the entire research pipeline for this paper was carried out by
Colombo, powered by GPT-5.6 Sol.
\par
\end{samepage}

\section{Introduction}\label{sec:intro}

Stochastic optimization guarantees depend on both the regularity of the
problem and the information available to the algorithm. A particularly
important distinction is whether a random sample may be evaluated more than
once. Under mean-squared smoothness, two gradients computed from the same
sample provide an accurate estimate of a gradient difference. If every
evaluation instead uses a fresh independent sample, that difference is not
directly available. This paper determines the minimax cost of this
single-query restriction for smooth nonconvex optimization.

We consider a differentiable, bounded-below objective
$F(x)=\E_\xi f(x;\xi)$ on $\R^d$, an initial point $x_0$, and the stationarity
criterion
\begin{equation}\label{eq:intro-target}
 \E\|\nabla F(\widehat x)\|\le\eps.
\end{equation}
The initial gap is at most $\Delta$, and the stochastic gradient
$g(x;\xi)=\nabla f(x;\xi)$ is unbiased with variance at most $\sigma^2$.
Mean-squared smoothness requires
\begin{equation}\label{eq:intro-mss}
 \E_\xi\|g(x;\xi)-g(y;\xi)\|^2
 \le \bar L^2\|x-y\|^2
 \qquad\text{for all }x,y\in\R^d.
\end{equation}
The same latent sample $\xi$ appears in both terms. Jensen's inequality
implies that $F$ is $\bar L$-smooth, but~\eqref{eq:intro-mss} places additional
structure on the stochastic gradients themselves.

\paragraph{The role of sample access.}
At each call, our oracle returns one stochastic gradient at the selected
point and then discards the sample. This is the $K=1$ model. In the $K=2$
model, the algorithm may evaluate a fresh sample at two points selected
before observing that sample's replies. We count individual gradient
evaluations in both cases. SPIDER~\cite{spider} uses $K=2$ access and attains
cubic dependence on $1/\eps$ when the other positive parameters are fixed.
Arjevani et al.~\cite[Theorem~3 and Proposition~1]{arjevani} prove a matching
cubic lower bound, including for gradients of scalar sample losses.
Their Section~6 asks whether the fourth-order SGD rate is optimal when
mean-squared smoothness is retained but access is restricted to $K=1$.

\paragraph{Main result.}
We answer this question by proving
\begin{equation}\label{eq:intro-rate}
 \mathsf T_1^\star(\Delta,\bar L,\sigma,\eps)
 =\Theta\!\left(
   \frac{\Delta\bar L}{\eps^2}
   +\frac{\Delta\bar L\sigma^2}{\eps^4}
 \right),
 \qquad 0<\eps\le c_0\sqrt{\Delta\bar L}.
\end{equation}
Here $\mathsf T_1^\star$ is the smallest deterministic query budget that
guarantees~\eqref{eq:intro-target} uniformly over finite dimensions, and
$c_0>0$ is universal. The new contribution is the lower bound. The upper
bound follows from classical stochastic gradient descent~\cite{ghadimi}.
Our lower bound has three features:
\begin{itemize}
 \item It applies to arbitrary adaptive randomized algorithms, with
 unrestricted query norms, memory, and computation. The output need not be a
 previously queried point.
 \item It holds in the smaller class of scalar sample losses with globally
 $\bar L$-Lipschitz gradients, and therefore also in the mean-squared-smooth
 class.
 \item It includes every variance level $\sigma\ge0$ in the stated accuracy
 regime, recovering the deterministic rate at $\sigma=0$ and the stochastic
 rate when $\sigma$ is larger than $\eps$.
\end{itemize}
The dimension of the hard instance may depend on the query budget. The
theorem is therefore a dimension-uniform characterization; it does not
assert the same rate in each fixed dimension.

\begin{table}[t]
\centering
\small
\renewcommand{\arraystretch}{1.35}
\setlength{\tabcolsep}{4pt}
\caption{Complexity for $\E\|\nabla F(\widehat x)\|\le\eps$, counting
individual gradient evaluations. All rows allow dimensions to grow and use
a variance bound $\sigma^2$. MSS denotes~\eqref{eq:intro-mss}; samplewise
smoothness means that every sample gradient is $\bar L$-Lipschitz. All rows
hold for $0<\eps\le c\sqrt{\Delta\bar L}$ and $\sigma\ge0$, with universal
constants that may differ. The earlier lower bound includes the
deterministic baseline; the last row combines SPIDER with SGD in the
low-noise regime, as explained in the text.}
\label{tab:comparison}
\begin{tabularx}{\textwidth}{@{}>{\raggedright\arraybackslash}p{0.255\textwidth}
 >{\centering\arraybackslash}p{0.065\textwidth}
 >{\centering\arraybackslash}p{0.36\textwidth}
 >{\raggedright\arraybackslash}X@{}}
\toprule
Result and regularity & Access & Query bound & Accuracy regime\\
\midrule
SGD upper bound~\cite{ghadimi}\newline
$\bar L$-smooth population
 & $K=1$
 & $\displaystyle O\!\left(\frac{\Delta\bar L}{\eps^2}
             +\frac{\Delta\bar L\sigma^2}{\eps^4}\right)$
 & $0<\eps\le c\sqrt{\Delta\bar L}$; $\sigma\ge0$\\[4pt]
Earlier lower bounds~\cite{arjevani,carmon}\newline
MSS; scalar sample gradients
 & $K=1$
 & $\displaystyle\Omega\!\left(\frac{\Delta\bar L}{\eps^2}
                   +\frac{\Delta\bar L\sigma}{\eps^3}
                   +\frac{\sigma^2}{\eps^2}\right)$
 & $0<\eps\le c\sqrt{\Delta\bar L}$; $\sigma\ge0$\\[4pt]
\textbf{This work}\newline
MSS; lower bound already under samplewise smoothness
 & $K=1$
 & $\displaystyle\Theta\!\left(\frac{\Delta\bar L}{\eps^2}
                  +\frac{\Delta\bar L\sigma^2}{\eps^4}\right)$
 & $0<\eps\le c_0\sqrt{\Delta\bar L}$; $\sigma\ge0$\\[4pt]
SPIDER/SGD and matching lower bounds~\cite{spider,ghadimi,arjevani,carmon}\newline
MSS; scalar sample gradients
 & $K=2$
 & $\displaystyle\Theta\!\left(\frac{\Delta\bar L}{\eps^2}
                   +\frac{\Delta\bar L\sigma}{\eps^3}
                   +\frac{\sigma^2}{\eps^2}\right)$
 & $0<\eps\le c\sqrt{\Delta\bar L}$; $\sigma\ge0$\\
\bottomrule
\end{tabularx}
\end{table}

The full-parameter comparisons in~\Cref{tab:comparison} combine existing
bounds rather than quoting a single theorem verbatim. Adding the
deterministic lower bound changes the maximum of the deterministic and
stochastic terms by at most a constant factor. For $\sigma\le\eps$, SGD
achieves $O(\Delta\bar L/\eps^2)$ in the stated regime. For
$\sigma\ge\eps$, the SPIDER rate
$O(\Delta\bar L\sigma/\eps^3+\sigma^2/\eps^2)$ also dominates the
deterministic term. This gives the uniform $K=2$ expression in the table.

\paragraph{Why the lower bound requires a different construction.}
The stochastic cost must remain large while the sample gradients satisfy
the common-sample constraint~\eqref{eq:intro-mss}. We ensure this regularity
by adding a linear Gaussian perturbation to a deterministic potential:
\[
 f_U(X,z;\xi)=F_U(X,z)+\langle\xi,z\rangle.
\]
The perturbation contributes the same vector at every evaluation point, so
it cancels from every same-sample gradient difference. Thus the sample
smoothness constant equals the population smoothness constant. All
statistical difficulty must be encoded in the geometry of $F_U$.

The potential consists of a sequence of gates associated with hidden
orthogonal directions. A point with small gradient must complete every
gate. Each unfinished gate exposes its direction through the same noisy
variable block, keeping the total variance independent of the number of
stages. The main complication is that a scalar potential also produces
exact gradients in the other blocks. By changing previously visited
coordinates, an algorithm can use these gradients to obtain information
about the current direction. The gates are designed so that, before a
stage is discovered, all such information reduces to one scalar projection
per response. A spherical information bound controls these exact
projections together with the noisy vector observation.

The resulting construction has
$\Theta(\Delta\bar L/\eps^2)$ stages, each requiring
$\Omega(1+\sigma^2/\eps^2)$ responses. A deferred-decision interaction
keeps future directions unexposed, and disjoint response counts allow the
stage costs to be added. A bounded radial encoding and a coercive guard
make the argument valid for arbitrarily large query and output norms.

\paragraph{Organization.}
\Cref{sec:related} places the result in the literature.
\Cref{sec:model,sec:upper} formalize the model and prove the upper bound.
\Cref{sec:construction} constructs the hard sample losses and establishes
their regularity and gradient barrier. \Cref{sec:information,sec:lower}
prove the information bounds and combine them into the lower bound.
\Cref{sec:discussion} discusses the scope of the characterization.
\Cref{app:probability,app:yao} collect the probability estimates and the
fixed-instance randomized conclusion.

\section{Related Work}\label{sec:related}

\paragraph{Stochastic descent.}
Ghadimi and Lan~\cite{ghadimi} analyze stochastic gradient methods for
nonconvex objectives and obtain stationarity guarantees by returning a
randomly selected iterate. We include the corresponding fixed-budget
argument in~\Cref{sec:upper}, so that the output rule and the expected-norm
criterion agree exactly with our minimax formulation. This upper bound
requires only population smoothness and bounded variance.

\paragraph{Variance reduction and common-sample differences.}
SPIDER~\cite{spider} and SNVRG~\cite{snvrg} improve the online accuracy
dependence using recursive or nested gradient estimators. STORM~\cite{storm}
and PAGE~\cite{page} provide further estimator designs. Their correction
steps evaluate the same sample at two points. Such corrections are
compatible with $K=2$ access, whereas a $K=1$ oracle supplies an independent
sample for each evaluation. The comparison in~\Cref{tab:comparison} keeps
this distinction explicit.

Fradin et al.~\cite{fradin2026} study bounded central moments and weak average
smoothness. Their $p=q=2$ specialization recovers bounded variance and
mean-squared smoothness, with cubic accuracy dependence using
common-sample differences. Their Algorithm~1 uses two evaluations of each
sample in its recursive correction, and their general-moment lower bounds
are formulated for zero-respecting algorithms. These results concern a
different access model from the unrestricted single-query problem studied
here.

\paragraph{Lower bounds for stationarity.}
Carmon et al.~\cite{carmon} develop smooth chain constructions and random
embeddings for stationary-point lower bounds; their first-order
specialization supplies the deterministic
$\Omega(\Delta\bar L/\eps^2)$ benchmark. Arjevani et
al.~\cite{arjevani} extend stochastic lower bounds to unrestricted
randomized algorithms and scalar sample gradients. Our proof also hides
successive directions, while its additive noise and exact-projection
analysis preserve samplewise smoothness at the fourth-order single-query
rate.

For comparisons with oracle conventions in this literature, a lower bound
that permits function values also applies when only gradients are returned.
A requirement that the output be queried can be met by one additional
call at an otherwise unqueried output. Moreover, for fixed $K\in\{1,2\}$,
counting rounds or individual gradient evaluations changes complexity by
only a universal factor. These observations align the cited lower bounds
with the conventions in~\Cref{sec:model}.

\paragraph{Recent results under different oracle assumptions.}
Peng, Tang, and Jia~\cite{peng2026} establish an
$\Omega(n+\sqrt n\,\Delta L_{\max}/\eps^2)$ lower bound for
unrestricted randomized finite-sum algorithms under individual smoothness,
for $n$ above a universal threshold and
$\eps^2\le c\Delta L_{\max}$. Their incremental oracle lets an algorithm
select and revisit a component index. Our oracle instead draws a fresh
sample, and our class specifies a uniform variance budget $\sigma^2$.
The finite-sum rate therefore does not directly identify the complexity
studied here.

Jin~\cite{jin2026} proves the fourth-order single-query lower bound under
an almost-sure bound on the gradient error, in the regime
$\eps\le c\min\{\sigma,\sqrt{L\Delta}\}$. That oracle class allows
discontinuous and nonconservative replies, as stated in Remark~3.2 of that
paper. Our hard oracle has conservative, globally Lipschitz sample
gradients, but its Gaussian error has unbounded support and bounded
variance. The two results thus establish the same accuracy exponent under
different structural restrictions; neither lower-bound construction
provides the combined restriction of samplewise smoothness and uniformly
bounded noise.

\paragraph{Additional derivative information.}
Arjevani et al.~\cite{secondorder} show that stochastic Hessian-vector
information can support cubic accuracy complexity with fresh-sample
access under appropriate higher-order regularity and noise assumptions.
Their result identifies another way to obtain useful information about
gradient changes. Our characterization concerns first-order access under
mean-squared smoothness, with a lower bound that already holds for
samplewise Lipschitz gradients.

\section{Model and Main Result}\label{sec:model}

Throughout, $\|\cdot\|$ is the Euclidean norm and $\|\cdot\|_{\rm op}$ is
the associated operator norm. A function is smooth when its gradient is
globally Lipschitz. Constants denoted by $c,C,K$ are positive and universal
unless stated otherwise; their values may change from one occurrence to
the next.

\subsection{Problem class and oracle access}
Fix $\Delta,\Lb,\eps>0$ and $\sigma\ge0$. For a finite dimension $d$, let
\[
 F(x)=\E_{\xi\sim P} f(x;\xi),\qquad x\in\R^d,
\]
where each sample loss is differentiable, the displayed expectations are
well-defined, and $g(x;\xi):=\nabla_x f(x;\xi)$. The problem class consists
of all such triples $(F,f,P)$, in every finite dimension, satisfying the
following conditions for all $x,y\in\R^d$:
\begin{align}
 \E_\xi g(x;\xi)&=\nabla F(x),\label{eq:unbiased}\\
 \E_\xi\norm{g(x;\xi)-\nabla F(x)}^2&\le\sigma^2,\label{eq:variance}\\
 \E_\xi\norm{g(x;\xi)-g(y;\xi)}^2&\le\Lb^2\norm{x-y}^2
 \qquad\forall x,y.\label{eq:mss}
\end{align}
The population objective $F$ is bounded below, and the supplied initial
point $x_0$ obeys
\begin{equation}\label{eq:initial-gap}
 F(x_0)-\inf_xF(x)\le\Delta.
\end{equation}
Condition~\eqref{eq:mss} is mean-squared smoothness with a common latent
sample. Jensen's inequality gives
$\|\nabla F(x)-\nabla F(y)\|\le\Lb\|x-y\|$. Only the population objective
is required to be bounded below. The oracle access rule is specified
separately from these function-class assumptions.

\begin{definition}[Fresh single-query oracle]\label[definition]{def:oracle}
At call $t$, after observing its internal random seed and the preceding
query--response transcript, the algorithm selects an arbitrary measurable
query $x_t\in\R^d$.  The oracle draws a fresh independent $\xi_t\sim P$ and
returns only $g(x_t;\xi_t)$.  The sample handle is then discarded.  In
particular, the algorithm cannot request $g(y;\xi_t)$ at another point, choose
a sample index, inspect $f(\cdot;\xi_t)$, or obtain a Hessian-vector product.
This is the $K=1$ model.
\end{definition}

The dimension, initial point, and bounds $\Delta,\Lb,\sigma$, together with
$\eps$, are public inputs. Algorithms may use arbitrary measurable
adaptive query rules and unlimited computation and memory. Query norms
are unrestricted, and the final output need not be a queried point. A
query budget is a deterministic cap that must hold on every realization.

\begin{definition}[Dimension-uniform minimax complexity]\label[definition]{def:complexity}
Let $\Tstar(\Delta,\Lb,\sigma,\eps)$ be the least integer $N$ for which a
single algorithmic rule, accepting the dimension as an input and using at most $N$ calls on every run to the oracle in
\Cref{def:oracle} satisfies
\[
 \sup_{d<\infty}\ \sup_{(F,f,P)\text{ admissible in }\R^d}
 \E\norm{\nabla F(\widehat x)}\le\eps.
\]
The expectation includes the algorithm's seed and all oracle samples.  If no
such $N$ exists, the complexity is $+\infty$.
\end{definition}

\subsection{Minimax characterization}
\begin{theorem}[Sharp single-query complexity]\label[theorem]{thm:main}
There are universal constants $c,C,c_0>0$ such that, whenever
\[
 0<\eps\le c_0\sqrt{\Lb\Delta},
\]
one has
\begin{equation}\label{eq:main-rate}
 c\left(
   \frac{\Delta\Lb}{\eps^2}
   +\frac{\Delta\Lb\sigma^2}{\eps^4}
 \right)
 \le \Tstar(\Delta,\Lb,\sigma,\eps)
 \le
 C\left(
   \frac{\Delta\Lb}{\eps^2}
   +\frac{\Delta\Lb\sigma^2}{\eps^4}
 \right).
\end{equation}
The lower bound holds against arbitrary adaptive randomized algorithms.
For each proposed finite horizon, the hard distribution is supported in
one finite dimension. It remains valid when the class is restricted to
scalar sample losses satisfying, almost surely,
\[
 \|\nabla f(x;\xi)-\nabla f(y;\xi)\|
 \le\Lb\|x-y\|\qquad\text{for every }x,y.
\]
\end{theorem}

\begin{corollary}[Precision exponent]\label[corollary]{cor:precision}
For fixed positive $\Delta,\Lb,\sigma$, as $\eps\downarrow0$ through the
nontrivial regime,
\[
 \Tstar(\Delta,\Lb,\sigma,\eps)=\Theta(\eps^{-4}).
\]
When $\sigma=0$, the same theorem reduces to the deterministic
$\Theta(\Delta\Lb/\eps^2)$ rate.
\end{corollary}

\begin{remark}[Scope]\label[remark]{rem:scope}
The characterization uses an expected gradient norm and a deterministic
query cap. Its lower-bound dimension may grow with the parameters and the
cap. The accuracy restriction isolates the nontrivial regime: smoothness
and the gap bound imply
$\|\nabla F(x_0)\|^2\le2\Lb\Delta$, so the supplied point is already
acceptable without an oracle call when $\eps\ge\sqrt{2\Lb\Delta}$.
The theorem does not specify the transition between this zero-query
regime and its stated accuracy range.
\end{remark}

\subsection{Proof strategy}
The upper bound is fixed-step stochastic gradient descent. The lower
bound combines an optimization barrier with a statistical cost. The hard
potential has
\[
 T=\Theta(\Delta\Lb/\eps^2)
\]
sequential stages. The construction in~\Cref{sec:construction} ensures
that every small-gradient output completes all of them, while preserving
a smoothness constant independent of $T$. The information analysis
in~\Cref{sec:information} shows that, before a stage is discovered, its
hidden direction can influence a reply only through one exact scalar
projection and a Gaussian vector with signal norm $O(\eps)$ and total
noise variance $\sigma^2$. This implies a response cost of
$\Omega(1+\sigma^2/\eps^2)$ for each stage. Finally, an auxiliary interaction
keeps the next direction conditionally uniform until its stage starts.
\Cref{sec:lower} assigns disjoint responses to stages and multiplies the
number of stages by their individual cost.

\section{The Upper Bound}\label{sec:upper}
The upper half of~\Cref{thm:main} is the standard SGD
guarantee~\cite{ghadimi}. The proof below records a deterministic query cap
and a directly implementable output rule.

\begin{lemma}[Fixed-budget SGD]\label[lemma]{lem:sgd}
In the accuracy regime of \Cref{thm:main}, there is a fresh-gradient algorithm with
\[
 N\le C\left(
  \frac{\Delta\Lb}{\eps^2}
  +\frac{\Delta\Lb\sigma^2}{\eps^4}
 \right)
\]
that returns \(\widehat x\) satisfying
\(\E\norm{\nabla F(\widehat x)}\le\eps\).
\end{lemma}
\begin{proof}
Condition~\eqref{eq:mss} and Jensen imply that \(F\) is \(\Lb\)-smooth.  Take
\[
 \eta=\frac1{2\Lb}\min\left\{1,\frac{\eps^2}{\sigma^2}\right\},
 \qquad
 N=\left\lceil\frac{4\Delta}{\eta\eps^2}\right\rceil,
\]
where the minimum is interpreted as $1$ when $\sigma=0$. Run
\(x_{t+1}=x_t-\eta g(x_t;\xi_t)\) and return an independent uniformly selected iterate from \(x_0,\ldots,x_{N-1}\).  Let $\mathcal F_t$ contain the internal seed and all samples before call $t$, so $x_t$ is $\mathcal F_t$-measurable and $\xi_t$ is independent of $\mathcal F_t$. Smoothness and conditional unbiasedness give
\[
 \E[F(x_{t+1})\mid\mathcal F_t]
 \le F(x_t)-\frac\eta2\norm{\nabla F(x_t)}^2
       +\frac{\Lb\eta^2\sigma^2}{2}.
\]
Telescoping yields
\[
 \frac1N\sum_{t<N}\E\norm{\nabla F(x_t)}^2
 \le\frac{2\Delta}{\eta N}+\Lb\eta\sigma^2\le\eps^2.
\]
The random output rule makes the left-hand side equal to
$\E\|\nabla F(\widehat x)\|^2$. Jensen's inequality gives the
expected-norm guarantee. Finally,
\[
 N\le 1+\frac{8\Delta\Lb}{\eps^2}
                 \max\left\{1,\frac{\sigma^2}{\eps^2}\right\},
\]
and the additive $1$ is absorbed by the claimed bound in the accuracy
regime of~\Cref{thm:main}.
\end{proof}

\section{A Smooth Hard Family}\label{sec:construction}
The hard family must satisfy three requirements at once. Its sample gradients
must be globally Lipschitz; an unfinished stage must keep the population
gradient bounded away from zero; and the oracle must reveal little about the
direction needed to complete that stage. We establish the first two
properties in this section and analyze the third in
\Cref{sec:information}.

The construction has two components. A bounded radial map limits the effect
of queries with large norm, while a radial penalty prevents those queries
from producing small gradients. A sequence of scalar gates then makes
stationarity require completion of every stage. The derivative bounds for
this sequence are independent of its length, which is essential for the
final scaling. Throughout, smoothness means a globally Lipschitz gradient:
the scalar switches are $C^{1,1}$, and Hessian identities are used almost
everywhere. No $C^2$ regularity is required.

\subsection{A bounded radial encoding}
For \(w\in\R^D\), define
\begin{equation}\label{eq:rho}
 \rho(w)=\frac{w}{\sqrt{1+\norm w^2}},
 \qquad J(w)=D\rho(w)
 =aI-a^3ww^\top,
 \quad a=(1+\norm w^2)^{-1/2}.
\end{equation}
The map and its first two derivatives satisfy the dimension-independent
bounds
\begin{equation}\label{eq:rho-bounds}
 \norm{\rho(w)}<1,
 \qquad \norm{J(w)}_{\rm op}\le1,
 \qquad \norm{D^2\rho(w)}_{\rm op}\le K
\end{equation}
for a universal \(K\).  The scalar guard
\begin{equation}\label{eq:guard}
 R(w)=\sqrt{1+\norm w^2}-1
\end{equation}
is nonnegative and coercive, with $\nabla R(w)=\rho(w)$ and
$\nabla^2R(w)=J(w)$. In particular, its gradient is globally Lipschitz.

\subsection{A radial gradient barrier}
The following lemma identifies the only possible cancellation between the
encoded hidden direction and the radial penalty. Such cancellation requires
a nontrivial correlation with the hidden direction.

\begin{lemma}[Radial hidden-direction coercivity]\label[lemma]{lem:radial-coercive}
Fix \(0<a_0\le a_1<\infty\) and \(\gamma>0\).  There are constants
\(q_*,\delta_*>0\), depending only on \((a_0,a_1,\gamma)\), with the following property.  For every unit vector \(u\), every \(w\in\R^D\), and every \(a\in[a_0,a_1]\),
\begin{align}
 \ip{u}{\rho(w)}\le q_*/2
 &\quad\Longrightarrow\quad
 \norm{-aJ(w)u+\gamma\rho(w)}\ge\delta_*,\label{eq:x-coercive}\\
 \ip{u}{\rho(w)}\ge-q_*/2
 &\quad\Longrightarrow\quad
 \norm{ aJ(w)u+\gamma\rho(w)}\ge\delta_*.
 \label{eq:z-coercive}
\end{align}
\end{lemma}
\begin{proof}
We first locate the zeros of the first vector field. Write
$w=tu+v$, where $v\perp u$. Its component orthogonal to $u$ is
\[
 a_w\bigl(\gamma+a a_w^2t\bigr)v,
 \qquad a_w=(1+\norm w^2)^{-1/2}.
\]
For $t\ge0$, the coefficient is positive, so a zero requires $v=0$.
For $t<0$, both terms in the component along $u$ are strictly negative.
The case $t=0$ also cannot give a zero. Thus every zero has
$w=tu$ with $t>0$. With
\(q=t/\sqrt{1+t^2}\), the zero equation is
\begin{equation}\label{eq:radial-root}
 a(1-q^2)^{3/2}=\gamma q.
\end{equation}
The left-hand side decreases strictly with $q\in(0,1)$, while the
right-hand side increases strictly. Equation~\eqref{eq:radial-root}
therefore has a unique solution $q_a\in(0,1)$, which increases with $a$.
Set $q_*=\min_{a\in[a_0,a_1]}q_a>0$. The zeros of the second vector
field are obtained by reflection, and have encoded correlation $-q_a$.

We next obtain a lower bound on the norm that is uniform in the ambient
dimension. Each norm depends only on the three scalar variables
$a,t,\|v\|$. On a bounded set of these variables, continuity and the
absence of zeros in the stated regions give a positive minimum. Moreover,
$\|J(w)\|_{\rm op}\le(1+\|w\|^2)^{-1/2}\to0$ and
$\|\rho(w)\|\to1$ as $\|w\|\to\infty$. Both vector-field norms
therefore tend to $\gamma$, uniformly for $a\in[a_0,a_1]$.
Combining the bounded and unbounded regions proves
\eqref{eq:x-coercive}--\eqref{eq:z-coercive}.
\end{proof}

\subsection{The smooth sequential latch}
We now choose the scalar coefficients to which the radial lemma will be
applied. Each stage has a completion coordinate $q_j$ and a readout
coordinate $s_j$. The completion gates select the first unfinished stage;
the readout at that stage determines whether progress can be made in
$q_j$. Flat portions of the gates will later ensure that undiscovered
directions remain absent from the oracle response.

All scalar functions introduced below are fixed independently of the stage
count and the optimization parameters. Only their stated bounds enter the
proof.

Choose once and for all constants
\[
 0<d_0<d_1<d_2<\tfrac14,
 \qquad 0<h_0<h_1<\tfrac14.
\]
Let the basic switch be
\[
 \vartheta(t)=
 \begin{cases}
 0,&t\le0,\\
 3t^2-2t^3,&0<t<1,\\
 1,&t\ge1.
 \end{cases}
\]
This nondecreasing $C^{1,1}$ switch satisfies, for a universal $K$,
\begin{equation}\label{eq:switch-energy}
 (\vartheta')^2\le K\vartheta(1-\vartheta),
 \qquad |\vartheta''|\le K\quad\text{a.e.}
\end{equation}
Indeed, on $(0,1)$,
\[
 \vartheta=t^2(3-2t),\qquad
 1-\vartheta=(1-t)^2(1+2t),\qquad
 \vartheta'=6t(1-t),
\]
so $(\vartheta')^2/[\vartheta(1-\vartheta)]
=36/[(3-2t)(1+2t)]\le36$, while $|\vartheta''|\le6$.
Define
\begin{align}
 C(q)&=\vartheta\!\left(\frac{q-d_1}{d_2-d_1}\right),
 &B(q)&=1-C(q),\label{eq:C-gate}\\
 A(s)&=1-\vartheta\!\left(\frac{s+h_1}{h_1-h_0}\right).
 \label{eq:A-gate}
\end{align}
Thus a stage is fully open when $q\ge d_2$ and fully closed when
$q\le d_1$. The readout switch satisfies $A=1$ on
$(-\infty,-h_1]$ and $A=0$ on $[-h_0,\infty)$. Absorbing the fixed
transition widths into $K$ gives
\begin{equation}\label{eq:C-energy}
 (C')^2\le KCB,
 \qquad |C''|+|A'|+|A''|\le K\quad\text{a.e.}
\end{equation}

Set $Q=2$ and $\eta=1$.  Choose
\[
 \alpha>1+\eta\|A'\|_\infty(Q+1),
 \qquad
 M_0>1+\alpha+\eta(Q+1),
\]
and put
\begin{equation}\label{eq:L-def}
 L(q,s)=M_0+\alpha s+\eta A(s)(Q-q).
\end{equation}
The choice of $\alpha$ ensures that the payoff increases with $s$, even
through the transition of $A$, while the choice of $M_0$ makes the payoff
positive. Specifically, on $[-1,1]^2$,
\begin{equation}\label{eq:L-properties}
 L\ge1,
 \qquad 1\le L_s(q,s)\le C_z,
 \qquad L_q(q,s)=-\eta A(s)\le0
\end{equation}
for a universal $C_z<\infty$.

The remaining local potential maintains a descent direction while a
completion gate is in transition. Define $V$ by $V(0)=0$ and
\begin{equation}\label{eq:V-properties}
 V'(q)=-\nu\,
 \vartheta\!\left(\frac{q-d_0}{d_1-d_0}\right).
\end{equation}
Here $\nu>0$ is a universal constant chosen after $C,A,L$ below.
Hence $V'=0$ on $(-\infty,d_0]$, $V'=-\nu$ on
$[d_1,\infty)$, $-\nu\le V'\le0$, and $V'$ is globally Lipschitz.

For \(q,s\in[-1,1]^T\), set
\begin{equation}\label{eq:prefixes}
 P_0=1,
 \qquad P_j=\prod_{i=1}^jC(q_i),
 \qquad w_j=P_{j-1}B(q_j),
\end{equation}
and define
\begin{equation}\label{eq:core}
 \mathcal J_T(q,s)
 =\sum_{j=1}^TV(q_j)
  +\sum_{j=1}^T w_jL(q_j,s_j).
\end{equation}
The weight $w_j$ represents the probability of reaching stage $j$ and
stopping there in a sequence of gates with continuation probabilities
$C(q_i)$. These weights telescope:
\begin{equation}\label{eq:hazard}
 \sum_{j=1}^Tw_j=1-P_T\le1.
\end{equation}

\begin{lemma}[Uniform calculus of the latch]\label[lemma]{lem:latch-calculus}
The fixed scalar functions above satisfy, uniformly over \(T\),
\begin{enumerate}
 \item \(|\mathcal J_T(q,s)|\le K(T+1)\);
 \item every coordinate of \(\nabla_q\mathcal J_T\) is bounded by \(K\),
       \(\norm{\nabla_s\mathcal J_T}_1\le K\), and
       \(\norm{\nabla^2\mathcal J_T}_{\rm op}\le K\) almost everywhere;
 \item if \(j\) is the first index with \(q_j<d_2\), then all earlier gates satisfy \(C(q_i)=1\), and:
 \begin{enumerate}
  \item if \(q_j<d_1\), then \(C(q_j)=C'(q_j)=0\), the only nonzero hidden-readout coefficient among stages \(j,j+1,\ldots\) is
  \[
    c_j:=\partial_{s_j}\mathcal J_T=L_s(q_j,s_j)\in[1,C_z],
  \]
  and
  \[
    b_j:=\partial_{q_j}\mathcal J_T=V'(q_j)-\eta A(s_j);
  \]
  \item if \(d_1\le q_j<d_2\), then \(b_j\le-b_0\) for a universal constant \(b_0>0\).
 \end{enumerate}
\end{enumerate}
\end{lemma}
\begin{proof}
Separate the weighted payoff from the local potentials:
\[
 H(q,s)=\sum_{j=1}^Tw_jL(q_j,s_j).
\]
Boundedness of $V$ and $L$ on their coordinate domains, together with
\eqref{eq:hazard}, gives the value bound. To control derivatives without
introducing a factor of $T$, interpret $H$ as the expectation of a bounded
payoff. Add an outcome $\infty$, corresponding to passage through all gates,
and define
\[
 \Pp(J=j)=w_j,
 \qquad
 \Pp(J=\infty)=P_T,
 \qquad L_\infty=0.
\]
Then $H=\E L_J$.

Fix a direction $h\in\R^T$ in the $q$ coordinates. Conditional on reaching
stage $i$, the experiment continues with probability $C_i=C(q_i)$ and stops
with probability $B_i=1-C_i$. The corresponding directional derivative of
the log transition probability is
\[
 \Delta_i(h)=
 \begin{cases}
 (C_i'/C_i)h_i,&\text{on continuation},\\
 -(C_i'/B_i)h_i,&\text{on stopping}.
 \end{cases}
\]
Branches of probability zero are omitted, and the realized increment is
zero at a flat endpoint. We first work away from the transition endpoints
and extend the resulting estimates below. The increments have conditional
mean zero and satisfy
\[
 \E[\Delta_i(h)^2\mid\text{stage $i$ is reached}]
 =\frac{(C_i')^2}{C_iB_i}h_i^2.
\]
Let $S_h$ be their sum up to the stopping outcome. Orthogonality of the
martingale differences gives
\begin{equation}\label{eq:fisher-hazard}
 \E S_h^2
 =\sum_{i=1}^TP_{i-1}
   \frac{(C_i')^2}{C_iB_i}h_i^2
 \le K\norm{h}^2.
\end{equation}

For a joint direction $d=(h,k)$ in $(q,s)$, differentiation of the finite sum
$H=\sum_j\Pp(J=j)L_j$ gives
\begin{equation}\label{eq:latch-first-directional}
 DH[d]=\E\bigl[DL_J[h_J,k_J]+L_JS_h\bigr].
\end{equation}
Because the first derivatives of $L$ are bounded,
\[
 \E|DL_J[h_J,k_J]|^2
 \le K\sum_jw_j(h_j^2+k_j^2)
 \le K\norm{d}^2.
\]
Equation~\eqref{eq:fisher-hazard} therefore bounds $\nabla H$ in Euclidean
norm. For the readout coordinates, the more explicit identity
\[
 \partial_{s_j}H=w_jL_s(q_j,s_j),
\]
and~\eqref{eq:hazard} give
$\|\nabla_sH\|_1\le K$.  The local term $\sum_jV(q_j)$ adds at most $\nu$ to
each $q$ coordinate, proving the stated first-derivative bounds.

Differentiating~\eqref{eq:latch-first-directional} once more in the same
direction yields
\begin{equation}\label{eq:latch-second-directional}
 D^2H[d,d]
 =\E\bigl[
 D^2L_J[d_J,d_J]+2DL_J[d_J]S_h
 +L_J(S_h^2+\dot S_h)
 \bigr],
\end{equation}
where $\dot S_h$ is the directional derivative of $S_h$. The terms involving
$D^2L_J$, $DL_J S_h$, and $L_JS_h^2$ have absolute expectation at most
$K\|d\|^2$, by the bounded payoff derivatives, Cauchy--Schwarz, and
\eqref{eq:fisher-hazard}. To bound the remaining term, observe that the two
possible second log transition derivatives at a reached stage are
\[
 \left(\frac{C_i''}{C_i}-\frac{(C_i')^2}{C_i^2}\right)h_i^2,
 \qquad
 \left(-\frac{C_i''}{B_i}-\frac{(C_i')^2}{B_i^2}\right)h_i^2.
\]
Consequently
\[
 \E|\dot S_h|
 \le\sum_iP_{i-1}
 \left(2|C_i''|+\frac{(C_i')^2}{C_i}
                    +\frac{(C_i')^2}{B_i}\right)h_i^2
 \le K\norm{h}^2
\]
by~\eqref{eq:C-energy}. Thus
$|D^2H[d,d]|\le K\|d\|^2$, and polarization gives
$\|\nabla^2H\|_{\rm op}\le K$. Adding the diagonal Hessian of
$\sum_jV(q_j)$ preserves a bound independent of $T$.

For each fixed $T$, the core is a finite sum of products of $C^{1,1}$
functions, so its gradient is locally Lipschitz with some finite bound.
The almost-everywhere estimate just proved makes this bound uniform in $T$:
mollify on the interior of the coordinate cube, integrate the bounded
Hessian along segments, and pass to the limit. Continuity extends the
result to the boundary. This also justifies the derivative estimates at
the flat transition endpoints.

Finally let $j$ be the first index with $q_j<d_2$.  Then $C(q_i)=1$ for every
$i<j$, so $P_{j-1}=1$.  If $q_j<d_1$, then $C(q_j)=C'(q_j)=0$ and every future
prefix vanishes.  Direct differentiation gives
\[
 c_j=L_s(q_j,s_j)\in[1,C_z],
 \qquad
 b_j=V'(q_j)-\eta A(s_j),
\]
while all later $c_i$ vanish.  If $d_1\le q_j<d_2$, write
\[
 R_{j+1}=\sum_{i>j}
 \left(\prod_{\ell=j+1}^{i-1}C(q_\ell)\right)
 B(q_i)L(q_i,s_i),
 \qquad |R_{j+1}|\le\sup|L|.
\]
Then
\[
 b_j=-\nu-C'(q_j)L(q_j,s_j)
     -B(q_j)\eta A(s_j)+C'(q_j)R_{j+1}.
\]
The other terms are uniformly bounded by constants that do not depend on
$\nu$. Choose, for example,
\[
 \nu>1+2\|C'\|_\infty\|L\|_\infty+\eta,
\]
where the norm of $L$ is over $[-1,1]^2$. This fixes $V$ and gives
$b_j\le-b_0$ with $b_0=1$.
\end{proof}

\subsection{Scalar sample losses and admissibility}
We now embed the scalar construction into Euclidean space. A separate block
$x_j$ records completion of stage $j$, while a shared block $z$ carries the
noisy readout used to unlock the current stage. Placing all noise in this
shared block keeps its total variance independent of the number of stages.
Fix integers \(T\ge1\) and \(D\ge T\).  Let
\(U=(u_1,\ldots,u_T)\in\R^{D\times T}\) be a Haar-uniform orthonormal frame.  The optimization variable is $v=(X,z)\in\R^{(T+1)D}$, where
\[
 X=(x_1,\ldots,x_T)\in(\R^D)^T,
 \qquad z\in\R^D.
\]
For a scale \(r>0\), define
\begin{equation}\label{eq:encoded-coordinates}
 q_j(X)=\ip{u_j}{\rho(x_j/r)},
 \qquad
 s_j(z)=\ip{u_j}{\rho(z/r)}.
\end{equation}
Thus \(|q_j|,|s_j|<1\).  Let \(\gamma_x,\gamma_z>0\) be fixed guards and define
\begin{equation}\label{eq:population-objective}
 F_U(X,z)
 =\lambda\mathcal J_T(q(X),s(z))
 +\gamma_x\lambda\sum_{j=1}^T R(x_j/r)
 +\gamma_z\lambda R(z/r).
\end{equation}
The initial point is the origin.

At each oracle call, draw an independent
$\xi\sim N(0,\sigma^2I_D/D)$ and return the gradient of
\begin{equation}\label{eq:sample-loss}
 f_U(X,z;\xi)=F_U(X,z)+\ip{\xi}{z}.
\end{equation}

Write $w_i=x_i/r$, $w_z=z/r$, and let
\[
 b_i=\partial_{q_i}\mathcal J_T(q,s),
 \qquad c_i=\partial_{s_i}\mathcal J_T(q,s),
 \qquad c=(c_1,\ldots,c_T).
\]
Since $J(w)=D\rho(w)$ is symmetric, the population-gradient blocks are
\begin{align}
 \nabla_{x_i}F_U(X,z)
 &=\frac{\lambda}{r}
   \bigl[b_iJ(w_i)u_i+\gamma_x\rho(w_i)\bigr],
 \label{eq:all-x-gradients}\\
 \nabla_zF_U(X,z)
 &=\frac{\lambda}{r}
   \bigl[J(w_z)Uc+\gamma_z\rho(w_z)\bigr].
 \label{eq:all-z-gradient}
\end{align}
The sample gradient has the same $x_i$ blocks; its $z$ block is
$\nabla_zF_U+\xi$.

\begin{lemma}[Admissibility and scaling]\label[lemma]{lem:admissible}
There is a universal \(K\) such that
\begin{align}
 \operatorname{Lip}(\nabla F_U)&\le K\lambda/r^2,\label{eq:smooth-scale}\\
 F_U(0)-\inf F_U&\le K\lambda(T+1).\label{eq:gap-scale}
\end{align}
For the oracle \eqref{eq:sample-loss},
\begin{align}
 \E\nabla f_U(X,z;\xi)&=\nabla F_U(X,z),\\
 \E\norm{\nabla f_U-\nabla F_U}^2&=\sigma^2,\\
 \E\norm{\nabla f_U(v;\xi)-\nabla f_U(v';\xi)}^2
 &=\norm{\nabla F_U(v)-\nabla F_U(v')}^2.
 \label{eq:samplewise-cancel}
\end{align}
Consequently, after choosing \(K\lambda/r^2\le\Lb\), every sample gradient is itself \(\Lb\)-Lipschitz and the instance satisfies \eqref{eq:unbiased}--\eqref{eq:mss}.
\end{lemma}
\begin{proof}
Let $\Phi(X,z)=(q(X),s(z))$. The $q$-Jacobian is block diagonal, with
operator norm at most $1/r$, while
\[
 Ds(z)=\frac1rU^\top J(z/r)
\]
has operator norm at most $1/r$. Since $q$ depends only on $X$ and $s$
only on $z$, the full Jacobian satisfies $\|D\Phi\|_{\rm op}\le1/r$.
The first term in the Hessian chain rule is
\[
 D\Phi^\top\nabla^2\mathcal J_TD\Phi,
\]
whose norm is at most $K/r^2$ by
\Cref{lem:latch-calculus}.

The remaining chain-rule terms account for the curvature of the coordinate
maps. The $q_i$ terms occupy disjoint $x_i$ blocks and have coefficients
$|b_i|\le K$, so
$\|\sum_i b_i\nabla^2q_i\|_{\rm op}\le K/r^2$.  In the shared block, hold the coefficient vector $c$ fixed at the point where the Hessian chain rule is evaluated. Then
\[
 \sum_i c_i\nabla^2s_i(z)
 =\nabla^2\ip{Uc}{\rho(z/r)},
\]
and its norm is at most
$K\|c\|_2/r^2\le K\|c\|_1/r^2\le K/r^2$.
The guard Hessians are bounded by a universal multiple of $\lambda/r^2$.
Multiplying the core bounds by $\lambda$ therefore gives the required
almost-everywhere Hessian estimate. The objective is locally $C^{1,1}$;
mollification on bounded open sets, followed by integration along segments
and passage to the limit, gives the global Lipschitz-gradient bound
\eqref{eq:smooth-scale}. This argument includes segments lying on a
switching hypersurface.

For the gap estimate, the core is bounded below by $-K(T+1)$: the weighted
payoff is nonnegative because
$w_j,L\ge0$, and each $V(q_j)$ is bounded on $[-1,1]$.  At the origin,
$V(0)=0$, $C(0)=0$, $A(0)=0$, and hence
$\mathcal J_T(0,0)=L(0,0)=M_0$.  The nonnegative guards are coercive in every
block, so $F_U$ is bounded below and
\[
 F_U(0)-\inf F_U\le K\lambda(T+1),
\]
which is~\eqref{eq:gap-scale}.

Finally, the sample perturbation is linear and centered.  Its gradient has
only a $z$-block, equal to $\xi$, so its second moment is
$\operatorname{tr}(\sigma^2I_D/D)=\sigma^2$. At two points evaluated with
the same latent sample, the vector $\xi$ cancels exactly. These observations
prove the three oracle identities. Equation~\eqref{eq:smooth-scale} then
gives the asserted samplewise $\Lb$-smoothness after rescaling.
\end{proof}

\begin{lemma}[Global stationarity barrier]\label[lemma]{lem:stationarity-barrier}
There are universal choices of the guard coefficients $\gamma_x,\gamma_z>0$
and a universal $\delta>0$ such that
\begin{equation}\label{eq:unscaled-barrier}
 \min_jq_j(X)<d_2
 \quad\Longrightarrow\quad
 \norm{\nabla F_U(X,z)}\ge2\delta\frac{\lambda}{r}
\end{equation}
for every frame and every $(X,z)$, including unbounded query or output norms.
Consequently, if $\lambda/r\ge4\eps/\delta$, then
\begin{equation}\label{eq:barrier}
 \min_j q_j(X)<d_2
 \quad\Longrightarrow\quad
 \norm{\nabla F_U(X,z)}>4\eps.
\end{equation}
\end{lemma}
\begin{proof}
The order of the constant choices is important. The scalar functions,
including $\nu$, are already fixed. Let $K_b$ and $b_0$ be the resulting
coefficient bounds from \Cref{lem:latch-calculus}, and let $C_z$ be the
constant in~\eqref{eq:L-properties}. We now choose the guards using
\Cref{lem:radial-coercive}.

First choose $\gamma_x>0$ sufficiently small that, for coefficients in
\[
 [a_x^-,a_x^+]=[\min\{b_0,\eta\},K_b],
\]
the root threshold $q_x$ satisfies $q_x>2d_2$.
Second choose $\gamma_z>0$ sufficiently small that, for coefficients in
$[1,C_z]$, the corresponding threshold $q_z$ satisfies $q_z>2h_1$.
These choices are possible because the positive root of
\eqref{eq:radial-root} tends to $1$ as $\gamma\downarrow0$, uniformly
over either fixed coefficient interval. Let $\delta>0$ be half the smaller
of the two resulting coercivity constants.

Let $j$ be the first index with $q_j<d_2$.  Every earlier gate is complete, so
$C(q_i)=1$ for $i<j$.

If $d_1\le q_j<d_2$, the completion gate is in transition and
$b_j\le-b_0$. Writing $a=-b_j$, the $x_j$ block is
\begin{equation}\label{eq:x-block-gradient}
 \nabla_{x_j}F_U
 =\frac\lambda r\left[-aJ(x_j/r)u_j
                 +\gamma_x\rho(x_j/r)\right],
 \qquad a\in[b_0,K_b].
\end{equation}
Because $q_j<d_2<q_x/2$, the first part of
\Cref{lem:radial-coercive} gives norm at least
$2\delta\lambda/r$.

If instead $q_j<d_1$, the completion gate is closed. Then $C(q_j)=0$,
all later prefixes vanish, and
all earlier readout coefficients vanish because their gates have $B(q_i)=0$.
Thus the shared block contains exactly one hidden direction:
\begin{equation}\label{eq:z-block-gradient}
 \nabla_zF_U
 =\frac\lambda r\left[c_jJ(z/r)u_j
                 +\gamma_z\rho(z/r)\right],
 \qquad c_j\in[1,C_z].
\end{equation}
If $s_j(z)\ge-q_z/2$, the second part of
\Cref{lem:radial-coercive} gives the desired lower bound in the $z$-block.
Otherwise $s_j<-q_z/2<-h_1$, so $A(s_j)=1$.  The latch identity then gives
\[
 b_j=V'(q_j)-\eta\le-\eta.
\]
With $a=-b_j\in[\eta,K_b]$ and
$q_j<d_1<d_2<q_x/2$, the first part of the radial lemma gives the same lower
bound in the $x_j$-block.

In every case, one gradient block has norm at least
$2\delta\lambda/r$, proving~\eqref{eq:unscaled-barrier}. Under the
stated scaling this lower bound is at least $8\eps$, so
\eqref{eq:barrier} holds with a strict margin. The radial lemma applies on
all of $\R^D$, and thus the same conclusion holds for queries and outputs
of arbitrarily large norm.
\end{proof}

\section{The Information Cost of a Stage}\label{sec:information}

The hard family exposes a hidden direction through two different channels.
The shared block returns a noisy vector, while exact gradients in earlier
blocks may reveal a noiseless scalar projection.  Both must be accounted for:
ignoring the exact channel would leave open a way to bypass the statistical
cost.  We first prove an information bound for adaptive access to both
channels.  We then show that the gradient oracle has precisely this form
until a stage is discovered, and formalize the sequential interaction needed
to apply the bound repeatedly.

\subsection{An adaptive spherical channel with exact projections}

\begin{lemma}[Adaptive posterior after exact linear measurements]\label[lemma]{lem:adaptive-posterior}
Let $W_0\subseteq\R^M$ have dimension $m$, and let $u$ be uniform on its unit
sphere. An adaptive procedure makes $K$ exact linear measurements, where $0\le K\le m/2$.  Before
measurement $\ell$, it chooses $a_\ell\in\R^M$ from the preceding transcript
and independent auxiliary randomness, and then observes $\ip{a_\ell}{u}$.
Additional observations are allowed between measurements, provided that their
conditional laws depend on $u$ only through the transcript and the exact
projections already revealed.  In particular, these observations contain no
information about the remaining orientation of $u$.

Let $S_K=\operatorname{span}\{P_{W_0}a_1,\ldots,P_{W_0}a_K\}$ and let
$r=\dim S_K$.  Conditional on the entire transcript,
\begin{equation}\label{eq:posterior-decomposition}
 u=p+\sqrt{1-\norm{p}^2}\,v,
 \qquad p=P_{S_K}u,
\end{equation}
where $v$ is uniform on the unit sphere of $W_0\cap S_K^\perp$.  Moreover,
for every $K\le m/2$ and every $\rho\in(0,1)$,
\begin{equation}\label{eq:adaptive-projection-tail}
 \Pp\!\left(\norm{P_{S_K}u}>\frac\rho2\right)
 \le 2\exp(-c\rho^2m)
 \qquad\text{whenever }K\le c\rho^2m,
\end{equation}
for a universal constant $c>0$.
\end{lemma}

\begin{proof}
\noindent\textbf{Conditional rotational invariance.}
Orthogonalize each selected vector against the preceding measurement span and
ignore a measurement whose residual component is zero.  Enumerate only the
rank-increasing measurements.  Normalize the new residual direction.  If $d$
is the current residual dimension, rotational invariance implies that the
newly observed normalized residual coordinate has the law of the first
coordinate of a uniform point on $\mathbb S^{d-1}$.  Conditional on that
coordinate, the remaining normalized residual is uniform on the corresponding
$(d-2)$-sphere.  Any extra observation allowed in the statement does not
change this conditional law.  Induction proves
\eqref{eq:posterior-decomposition}.

\noindent\textbf{The energy captured by adaptive measurements.}
Suppose the original interaction has only $r\le K$ rank-increasing
measurements.  After the interaction is over, append $K-r$ canonical exact
projections in the final residual subspace, without returning them to the
algorithm or changing any of its previous choices.  The enlarged measured
subspace contains the original one, so its projected energy dominates
$\norm{P_{S_K}u}^2$.  At each of the resulting $K$ effective rank increments,
whether it came from the algorithm or from this terminal padding, the squared
normalized residual coordinate has law
$\operatorname{Beta}(1/2,(d-1)/2)$ when the current residual dimension is $d$,
and this conditional law is independent of the preceding history.  These are
exactly the neutral stick-breaking ratios obtained from independent
$G_1^2,\ldots,G_m^2\sim\chi_1^2$:
\[
 Z_\ell^2=\frac{G_\ell^2}{G_\ell^2+\cdots+G_m^2},
 \qquad 1\le\ell\le K.
\]
Therefore the enlarged measured energy has the exact law
\begin{equation}\label{eq:adaptive-beta}
 \frac{X}{X+Y},
 \qquad X=\sum_{\ell\le K}G_\ell^2\sim\chi_K^2,
 \quad Y=\sum_{\ell>K}G_\ell^2\sim\chi_{m-K}^2,
\end{equation}
with $X,Y$ independent.  This is the same law as the energy in a fixed
$K$-dimensional subspace and yields an upper tail bound for the original,
possibly rank-deficient, adaptive measurement sequence.

For $K=0$, the tail bound is immediate.  Otherwise, assume
$K\le\rho^2m/128$.  The upper-tail estimate in \Cref{lem:chi-square} applies
at $\rho^2m/32\ge4K$.  For the independent variable $Y$, Chernoff's
inequality gives
\[
 \Pp(Y<m/2)\le e^{m/4}2^{-(m-K)/2}\le e^{-m/16}.
\]
Consequently,
\[
 \Pp(X>\rho^2m/32)+\Pp(Y<m/2)
 \le e^{-\rho^2m/256}+e^{-m/16}
 \le 2e^{-\rho^2m/256}.
\]
Outside these events, $X/(X+Y)\le\rho^2/16<\rho^2/4$.  Renaming constants
proves~\eqref{eq:adaptive-projection-tail}.
\end{proof}

\begin{lemma}[Adaptive exact-projection/Gaussian channel]\label[lemma]{lem:side-channel}
Let $u$ be uniform on the unit sphere of a fixed $m$-dimensional subspace of
$\R^M$, where $M\le2m$.  In each of $n$ adaptive rounds an algorithm may:
\begin{enumerate}
 \item obtain at most $k$ exact scalar projections of $u$, with the measurement
 vectors chosen adaptively; and
 \item observe a vector that, conditional on $u$ and the transcript so far, is Gaussian with covariance $\sigma^2 I_M/M$.
\end{enumerate}
Suppose that a reference experiment uses the same prior, exact-measurement
kernels, and adaptive query rules.  Its Gaussian mean must be determined by
the observed history and current exact projections, without further dependence
on $u$.  Assume that the true and reference Gaussian means differ in norm by
at most $a>0$ at every round and history.

The algorithm also selects at most $b(n+1)$ test vectors, collected in
$\mathcal V$, each of norm at most one; here $b$ is a fixed universal
constant.  Each vector is measurable with respect to the transcript and
independent auxiliary randomness.  Its true correlation with $u$, or the
outcome of a correlation test, is not supplied as an additional observation.
For every $\rho\in(0,1)$ there are constants $c,C>0$, depending only on $b$,
such that
\begin{equation}\label{eq:side-conditions}
 kn\le c\rho^2m,
 \qquad m\ge C\rho^{-2}\log(2n+2),
 \qquad n\le c\rho^2\frac{\sigma^2}{a^2}
\end{equation}
imply
\begin{equation}\label{eq:side-conclusion}
 \Pp\left(\max_{v\in\mathcal V}|\ip{u}{v}|\ge\rho\right)\le\frac18.
\end{equation}
When $\sigma=0$, the last condition is interpreted as $n=0$.
\end{lemma}

\begin{proof}
If $\sigma=0$, the stated convention forces $n=0$; the conclusion is then the
spherical-cap bound (and the rank condition has $K=0$).  Assume henceforth
that $\sigma>0$.

\noindent\textbf{Divergence between the two experiments.}
Let $P$ and $Q$ denote the joint laws of $u$ and the entire transcript in the
true and reference experiments.  We include the latent $u$ in both probability
spaces.  The prior and all exact-measurement kernels are identical.  At a fixed
history, the conditional KL divergence between the two Gaussian kernels is at
most
\[
 \frac{M}{2\sigma^2}\norm{\mu_t(u)-\bar\mu_t}^2
 \le \frac{Ma^2}{2\sigma^2}.
\]
The adaptive chain rule therefore yields
\begin{equation}\label{eq:side-kl}
 \KL(P\|Q)\le\frac{nMa^2}{2\sigma^2}
 \le\frac{nma^2}{\sigma^2}.
\end{equation}

\noindent\textbf{A rare event under the reference law.}
Under $Q$, the Gaussian observations provide no information about the remaining
orientation of $u$ beyond the exact projections.  Apply \Cref{lem:adaptive-posterior} after all
at most $K=kn$ exact measurements.  On an event of $Q$-probability at least
$1-2e^{-c\rho^2m}$, the measured component $p$ in
\eqref{eq:posterior-decomposition} has norm at most $\rho/2$.  Conditional on
the full reference transcript, all test vectors are fixed and the residual
$v$ is uniform on a sphere of dimension at least $m-K\ge m/2$.  Hence the
spherical-cap bound and a union bound give, for the event $A$ in
\eqref{eq:side-conclusion},
\begin{equation}\label{eq:reference-small-event}
 Q(A)\le 2e^{-c\rho^2m}+2b(n+1)e^{-c\rho^2m}
 \le \exp(-c'\rho^2m),
\end{equation}
after increasing the constant in the dimensional condition.

\noindent\textbf{Transfer to the true law.}
For $p=P(A)$ and $q=Q(A)$, binary data processing gives
\[
 \KL(P\|Q)\ge \operatorname{kl}(p\|q)
 \ge p\log(1/q)-\log2.
\]
If $p>1/8$, the right-hand side is at least $c''\rho^2m$ under the same
dimensional condition.  This contradicts~\eqref{eq:side-kl} when the last
inequality in~\eqref{eq:side-conditions} has a sufficiently small universal
constant.  Thus $p\le1/8$.
\end{proof}

\subsection{What the exact gradient reveals}

The following calculation isolates all dependence on an undiscovered direction.
It also explains why returning to earlier gates does not create independent
exact measurements of that direction within a single call.  Define the tail
payoff
\begin{equation}\label{eq:tail-payoff}
 R_{i+1}(q,s)
 =\sum_{\ell>i}\left(\prod_{h=i+1}^{\ell-1}C(q_h)\right)
 B(q_\ell)L(q_\ell,s_\ell).
\end{equation}
Direct differentiation of~\eqref{eq:core} gives
\begin{align}
 c_i:=\partial_{s_i}\mathcal J_T
 &=P_{i-1}B(q_i)L_s(q_i,s_i),\label{eq:ci-formula}\\
 b_i:=\partial_{q_i}\mathcal J_T
 &=V'(q_i)+P_{i-1}\bigl[
 B(q_i)L_q(q_i,s_i)
 +C'(q_i)(R_{i+1}-L(q_i,s_i))\bigr].
 \label{eq:bi-formula}
\end{align}

\begin{lemma}[Locality at the current and next stages]\label[lemma]{lem:front-locality}
Fix an index $j$ and a query.  Suppose every $i>j$ satisfies
\begin{equation}\label{eq:future-flat}
 q_i\le d_0,\qquad s_i\ge-h_0.
\end{equation}
Then the following hold.
\begin{enumerate}
 \item If also $q_j\le d_0$ and $s_j\ge-h_0$, the entire exact part of the
 gradient depends on the unknown $u_j$ only through the single scalar
 $s_j=\ip{u_j}{\rho(z/r)}$, except for the shared-block vector
 \begin{equation}\label{eq:current-vector-channel}
  \frac\lambda r\,c_jJ(z/r)u_j,
  \qquad |c_j|\le K.
 \end{equation}
 No direction $u_i$ with $i>j$ occurs in the gradient.
 \item If the query discovers stage $j$, meaning
 $q_j>d_0$ or $s_j<-h_0$, then the gradient can involve unknown frame columns
 only through $u_j$ and, when $j<T$, $u_{j+1}$; no column $u_i$ with
 $i>j+1$ occurs.
\end{enumerate}
These statements remain true when any subset of the earlier blocks is moved
back into a transition region.
\end{lemma}

\begin{proof}
In the first case, flatness gives
\[
 C(q_j)=C'(q_j)=V'(q_j)=A(s_j)=A'(s_j)=0.
\]
Equations~\eqref{eq:ci-formula}--\eqref{eq:bi-formula} yield
$b_j=0$ and $c_j=P_{j-1}\alpha$.  Since $C(q_j)=0$, every later prefix
$P_{i-1}$ vanishes.  Under~\eqref{eq:future-flat}, also $V'(q_i)=0$; hence
$b_i=c_i=0$ for $i>j$.

For $i<j$, the only term in the tail~\eqref{eq:tail-payoff} containing the
current hidden column is the payoff at the first flat stage $j$.  Before
discovery,
\[
 L(q_j,s_j)=M_0+\alpha s_j,
\]
so every earlier exact block gradient depends on $u_j$ only through the same
scalar $s_j$.  The $x_j$-block has no hidden component because $b_j=0$.  The
shared $z$-block is the sum of known earlier directions, the vector in
\eqref{eq:current-vector-channel}, the radial guard, and the additive Gaussian
noise.  This proves the first assertion, including simultaneous resets of any
number of old gates.

For the second assertion, the case $j=T$ has no later columns and is
immediate.  If $j<T$, the flat coordinate $q_{j+1}\le d_0$ has
$C(q_{j+1})=C'(q_{j+1})=0$.  Thus every prefix and every derivative reaching
past $j+1$ vanishes.  Local terms at $i>j+1$ also have zero derivative by
\eqref{eq:future-flat}.  Hence at most the two consecutive columns $u_j$ and
$u_{j+1}$ occur.
\end{proof}

\begin{lemma}[Stage discovery cost]\label[lemma]{lem:stage-cost}
Fix the scalar functions and let $\lambda/r=c_*\eps$, where $c_*>0$ is
universal.  At the start of stage $j$, suppose that, conditional on the entire
stage-start history and all earlier frame columns, $u_j$ is uniform on the
unit sphere of an orthogonal complement of dimension $m\ge D/2$.
Before discovery, all later columns are suppressed as in
\Cref{lem:sequential-coupling}.  A query selected while stage $j$ is current
discovers it when
\begin{equation}\label{eq:discover-events}
 q_j(X)>d_0
 \qquad\text{or}\qquad
 s_j(z)<-h_0.
\end{equation}
The algorithm receives every exact gradient block and the noisy shared block.
A stage may begin after a query has been selected but before its response is
returned; in this case, its first assigned response is the flat-extension
response specified in \Cref{lem:sequential-coupling}.

There are constants $c,C>0$ such that, with
\begin{equation}\label{eq:n0}
 n_0=1+\left\lfloor c\frac{\sigma^2}{\eps^2}\right\rfloor
\end{equation}
and $m\ge C(n_0+\log(N+2))$, the conditional probability of discovering the
stage after fewer than $n_0$ assigned responses is at most $1/8$.
The bound is uniform over the stage-start history and all admissible query
rules, including rules that use the earlier revealed directions.
\end{lemma}

\begin{proof}
\noindent\textbf{Reduction to one exact projection and one noisy vector.}
Condition on a stage-start history that includes all earlier frame columns.  Before discovery, \Cref{lem:front-locality} shows that the full
collection of exact $x$-block gradients depends on $u_j$ only through
\[
 s_j=\ip{u_j}{\rho(z/r)}.
\]
Even if many old blocks are simultaneously reset, they produce only known
functions of this same scalar.  We therefore reveal $s_j$ itself after every
response; this is a single exact linear measurement and can only strengthen
the algorithm.

After subtracting the known earlier-direction and radial-guard parts, the
noisy $z$-block has conditional form
\[
 Y_t=\frac{\lambda}{r}c_{j,t}J(z_t/r)u_j+\xi_t,
 \qquad
 \xi_t\sim N(0,\sigma^2I_D/D),
 \qquad |c_{j,t}|\le K.
\]
The reference experiment deletes the displayed hidden-direction term and
keeps the same query policy, exact projections, known gradient blocks, and
Gaussian noise.  Conditional on the exact projection and the preceding
reference transcript, its Gaussian mean contains no residual information
about $u_j$.  The true and reference means differ by at most
\[
 a\le K\lambda/r\le K'\eps.
\]

\noindent\textbf{A continuation that does not report discovery.}
We apply the channel bound to a fixed-length experiment, preserving the
spherical prior at stage start.  In this experiment, the current column always
uses its undiscovered-stage values:
\[
 C=C'=V'=A=A'=0,\qquad
 L=M_0+\alpha s_j,\qquad L_s=\alpha,\qquad L_q=0.
\]
Earlier columns retain their actual values.  In this experiment,
$c_{j,t}=\alpha P_{j-1}$ is uniformly bounded and every exact reply depends
on $u_j$ only through $s_j$, for every query point.

Run the original query rule while always supplying the ``not yet discovered''
state.  In particular, this continuation supplies neither the outcomes of the
discovery tests nor the free revelation of $u_j$, even if a test succeeds.
On histories that the original interaction would not reach, extend the rule
arbitrarily as a measurable function of the available transcript.  Before the
first successful test, the continued and original responses and state
registers agree.  Their first successful test therefore occurs at the same
selected query.  This comparison uses the original spherical prior at stage
start, without conditioning it on previous test failures.

If the stage begins at a half-step, also test the query that was already
selected before $u_j$ was sampled.  This only enlarges the event being bounded;
it does not cause a second stage advance at that query.  Its response is the
flat-extension response in both interactions until the first successful test.

\noindent\textbf{Application of the channel bound.}
The alternatives in~\eqref{eq:discover-events} are one-sided correlation
tests against $\rho(x_j/r)$ and $\rho(z/r)$, whose norms are at most one.
Either alternative implies an absolute correlation of at least
$\rho_0=\min\{d_0,h_0\}$.  Apply \Cref{lem:side-channel} to the continued
experiment with one exact projection per response, two test vectors per
selected query, ambient dimension $M=D$, and residual dimension $m\ge D/2$.
There are at most $n+1$ selected test points after $n=n_0-1$ responses;
pad a shorter continuation without returning the padding to the original
algorithm.  The assumption $m\ge C(n_0+\log(N+2))$ ensures the rank and
spherical-cap conditions.  Choosing the constant in~\eqref{eq:n0} sufficiently
small also ensures $n\le c\rho_0^2\sigma^2/a^2$.  The probability of any
successful test in this continuation is therefore at most $1/8$, which bounds
the probability of a discovery after fewer than $n_0$ assigned responses.

If the integer in~\eqref{eq:n0} is $n_0=1$, no response is needed in the
comparison experiment.  The conclusion then follows directly from the
spherical-cap bound for the initial query, including a query selected before a
half-step stage start.  This covers $\sigma=0$ as well.
\end{proof}

\subsection{A sequential interaction with deferred directions}

A discovery query may already contain information about the next direction in
its response.  We assign this response to the next stage and sample that
stage's direction only after the query has been fixed.  The following
interaction formalizes this convention on every trajectory, including those
where the eventual coupling to the true oracle fails.

\begin{lemma}[Sequential truncation coupling]\label[lemma]{lem:sequential-coupling}
Fix a horizon $N$ and a deterministic adaptive algorithm.  There is an
auxiliary interaction and an augmented copy of the algorithm that executes
the original query rule on the ordinary response component and ignores extra
revelations.  At each selected query, only the current stage is tested for
discovery.  A success advances the current index by one, and the next stage
is not tested again at the same query.  The following properties hold.
\begin{enumerate}
 \item After a stage is discovered, its hidden direction is revealed to the
 algorithm for free, after the discovery query is selected.  The algorithm may
 use all revealed directions and may reset arbitrary old blocks in later
 queries.
 \item Before a response, if $j$ is the current stage and the selected query
 does not discover it, the auxiliary gradient is evaluated from the same
 formulas as the true gradient with every hidden-core contribution from
 columns after $j$ suppressed; all radial-guard gradients, which are known from
 the query and contain no hidden frame, are retained.  If the
 query discovers $j<T$, the next column is sampled from the deferred Haar
 completion after the query has been selected, and its contribution is
 inserted through the \emph{flat extension}: irrespective of the query's
 actual next-stage correlations, the auxiliary response uses
 \[
 C=C'=V'=A=A'=0,\qquad
 L=M_0+\alpha s_{j+1},\quad L_s=\alpha,\quad L_q=0
 \]
 for that newly introduced column. More explicitly, compute the core-gradient coefficients by differentiating the latch formulas with these substituted scalar values and derivatives; a zero gate kills every later prefix. This retains the tail values that occur in derivatives of earlier gates. Retain all radial guards.  The same fresh additive
 Gaussian noise is used.  Thus, on every auxiliary trajectory, the first
 response involving $u_{j+1}$ depends on it only through the exact scalar
 $s_{j+1}$ and one Gaussian mean shift of norm $O(\lambda/r)$.  This response is
 assigned to the next stage.  If $j=T$, reveal the final column and return the full gradient from then on; all next-column clauses are omitted.
 \item With
 \begin{equation}\label{eq:dimension-choice}
 D\ge 2T+C\rho_0^{-2}
 \bigl(n_0+\log(24T(N+1))\bigr),
 \qquad \rho_0=\min\{d_0,h_0\},
 \end{equation}
 the original algorithm's queries and ordinary oracle-response transcript
 coincide with those of its augmented auxiliary copy, including the final
 output as a response-free query, with probability at least $11/12$.
\end{enumerate}
The stage-cost and counting bounds also hold for query rules that use the
extra revelations.  The transcript coupling is invoked for the augmented copy
that ignores them.  All conditional-Haar calculations and probability bounds
use the unstopped auxiliary law, before any restriction to the coupling event.
\end{lemma}

\begin{proof}
\noindent\textbf{Conditional law of the unused columns.}
Sample the orthonormal frame by deferred decisions.  When stage $j$ is
current, every preceding auxiliary response depends on the frame only through
its first $j$ columns.  This is true on every auxiliary trajectory.
Conditional on the exposed prefix and on the latent current column, the
remaining frame is Haar on the unused orthogonal complement.  If the query discovers $j$, the column $u_{j+1}$ is sampled only
after the query is fixed and enters that response solely through the flat
extension specified in the statement.  Consequently the single-stage channel
description of \Cref{lem:front-locality} holds on every auxiliary
trajectory, including trajectories on which the later coupling event fails.

\noindent\textbf{A coupling event under the auxiliary law.}
For each query, impose the good projections
\begin{equation}\label{eq:future-good}
 |\ip{u_i}{\rho(x_i/r)}|\le\rho_0,
 \qquad
 |\ip{u_i}{\rho(z/r)}|\le\rho_0
 \quad\text{for every }i>j.
\end{equation}
Impose the same tests on the final output.  Their outcomes are used only in
this analysis and are not returned to the algorithm.  Each vector has norm
at most one and is fixed before the corresponding future column is exposed.
Conditional on the current prefix, each future column is marginally uniform
on a sphere of dimension at least $D-T$.  The spherical-cap bound and a union
bound over at most $2T(N+1)$ tests show that some inequality
in~\eqref{eq:future-good} fails with probability at most $1/12$ under
\eqref{eq:dimension-choice}.  No conditioning on success of previous tests is
used in this calculation.

\noindent\textbf{Equality of the ordinary transcripts.}
On the good event, every later column lies in both flat regions.  If the current
query does not discover $j$, the first part of
\Cref{lem:front-locality} says that the true response contains no later
column.  If it does discover $j$, the second part says that at most the next
column $u_{j+1}$ occurs.  On the good event that next column is itself in the
undiscovered plateau, so its true derivatives are exactly the flat-extension
values.  Hence the true and truncated auxiliary responses are identical.  Using the same frame, noises, and algorithm seed, induction gives identity
of the original queries and ordinary responses with those of the augmented
copy until the first violation of~\eqref{eq:future-good}.  The extra revealed
directions are placed in a separate side-information register and are ignored
by this copy when it executes the original query rule.  \Cref{lem:stage-cost}
is uniform over the larger class of rules that do use this register, so the
revelation can safely be retained in the stage analysis.
\end{proof}

\section{The Lower Bound}\label{sec:lower}

We now combine the analytic and information bounds.  The normalization
produces $\Theta(\Delta\Lb/\eps^2)$ stages while keeping the signal in each
noisy response of order $\eps$.  The sequential interaction then forces the
cost $\Omega(1+\sigma^2/\eps^2)$ to be paid separately at each stage.

\paragraph{Choice of constants and scales.}
Choose the scalar constants and scales in the following order.  First
choose $d_0<d_1<d_2$ and $h_0<h_1$, then choose $C,A,L,V$ and the transition
slope $\nu$.  Their derivative bounds determine finite coefficient
intervals in \Cref{lem:stationarity-barrier}.  Finally choose the fixed
guard coefficients $\gamma_x,\gamma_z>0$ sufficiently small that the roots in
\Cref{lem:radial-coercive} lie above $2d_2$ and $2h_1$, respectively.
All these numbers are universal and are now frozen.  Let $K_0$ dominate the universal
smoothness and gap constants in \Cref{lem:latch-calculus,lem:admissible}, as well as $4/\delta$ from
\Cref{lem:stationarity-barrier}.  Choose constants $c_r,c_\lambda$ so that
\begin{equation}\label{eq:constant-hierarchy}
 \frac{c_\lambda}{c_r}\ge K_0,
 \qquad
 K_0\frac{c_\lambda}{c_r^2}\le1.
\end{equation}
For example, take $c_r\ge K_0^2$ and set $c_\lambda=K_0c_r$.
Put
\begin{equation}\label{eq:scales}
 r=c_r\frac{\eps}{\Lb},
 \qquad
 \lambda=c_\lambda\frac{\eps^2}{\Lb},
 \qquad
 T=\left\lfloor c_T\frac{\Delta\Lb}{\eps^2}\right\rfloor,
\end{equation}
where $c_T=(4K_0c_\lambda)^{-1}$. Choose the constant in \Cref{thm:main} so that $c_0^2\le c_T/2$. Writing $A_0=\Delta\Lb/\eps^2$, we then have $c_TA_0\ge2$ and
\[
 \tfrac12c_TA_0\le T\le c_TA_0,
 \qquad T+1\le2c_TA_0,
 \qquad K_0\lambda(T+1)\le\Delta/2.
\]
It follows that $T\ge1$, the initial gap is at most $\Delta$, every sample
gradient is $\Lb$-Lipschitz, and the stationarity barrier is greater than
$4\eps$.  Moreover, $\lambda/r=K_0\eps$, as required by
\Cref{lem:stage-cost}.

\begin{theorem}[Conservative $K=1$ lower bound]\label[theorem]{thm:lower}
Fix $\Delta,\Lb>0$, $\sigma\ge0$, and $0<\eps\le c_0\sqrt{\Delta\Lb}$. In the model of \Cref{def:oracle}, every adaptive randomized algorithm that guarantees $\E\|\nabla F(\widehat x)\|\le\eps$ on all instances satisfying \eqref{eq:unbiased}--\eqref{eq:initial-gap} requires
\begin{equation}\label{eq:lower-rate}
 \Omega\!\left(
   \frac{\Delta\Lb}{\eps^2}
   +\frac{\Delta\Lb\sigma^2}{\eps^4}
 \right)
\end{equation}
fresh single-point sample-gradient calls. More precisely, with the preceding scales, every deterministic cap $N<Tn_0/2$ admits a finite-dimensional fixed frame $U$ for which $\E\|\nabla F_U(\widehat x)\|>\eps$.
\end{theorem}

\begin{proof}
\noindent\textbf{A common hard distribution.}
First fix a deterministic algorithm and draw a Haar frame.  Let $n_0$ be
given by~\eqref{eq:n0}.  Choose a finite $D$ satisfying
\eqref{eq:dimension-choice}, enlarged if necessary for
\Cref{lem:side-channel}.  This choice depends only on the public parameters
and $N$, and therefore works for every query policy and every realization of
an internal random seed.  An algorithm that stops early can be padded to
exactly $N$ calls while retaining its stored output and ignoring all padding
responses.  Run the auxiliary interaction of
\Cref{lem:sequential-coupling}.

\noindent\textbf{Assigning each response to one stage.}
Append the final output as a response-free selected query with index $N+1$.
Let $\tau_j$ be the first query index that discovers stage $j$ while it is
current, and let $\tau_j=\infty$ if no such query is selected.  The state
advances by at most one stage at each selected query.

Stage $1$ begins before query $1$.  For $j\ge2$, stage $j$ begins at the
half-step after query $\tau_{j-1}$ is selected and $u_{j-1}$ is revealed,
but before that query's response is returned.  At this half-step, $u_j$ is
uniform in the unused orthogonal complement, conditional on all earlier
columns and the complete history.  Indeed, the query was fixed before any
observation involving $u_j$, and the auxiliary construction samples $u_j$
only afterward.  If $\tau_{j-1}=N+1$, no response remains and the new stage
cannot be discovered.

Set $L_j=\infty$ for a stage not discovered by the output.  For a discovered
stage, let $L_j$ be the number of responses assigned to it before its
discovery query is selected:
\begin{equation}\label{eq:stage-lengths}
 L_1=\tau_1-1,
 \qquad
 L_j=\tau_j-\tau_{j-1}\quad(2\le j\le T).
\end{equation}
The response to the discovery query of stage $j-1$ is thus charged to
stage $j$.  It has the flat-extension form used in
\Cref{lem:stage-cost}, even though its query was selected before stage $j$
began.  The already selected point may be included among that lemma's test
vectors, without advancing the state twice.  Distinct stages receive disjoint
responses.  If all stages are discovered by the final output, telescoping gives
\begin{equation}\label{eq:length-telescope}
 \sum_{j=1}^TL_j=\tau_T-1\le N.
\end{equation}
\noindent\textbf{Few stages can be discovered quickly.}
Call a stage \emph{fast} if it starts and $L_j<n_0$, and set its fast
indicator to zero if it never starts.  At its start, the unused orthogonal
complement has dimension $m\ge D-T\ge D/2$.  Conditional on the complete
stage-start history, \Cref{lem:stage-cost} gives
\begin{equation}\label{eq:fast-prob}
 \Pp(L_j<n_0\mid\text{stage-$j$ start history})\le\frac18.
\end{equation}
Therefore the expected number of fast stages is at most $T/8$.  If every stage
is discovered with $N<Tn_0/2$ oracle responses, then
\eqref{eq:length-telescope} forces at least $T/2$ stages to be fast.  Markov's
inequality yields
\begin{equation}\label{eq:all-stages-prob}
 \Pp(\text{all $T$ stages are discovered by the output})\le\frac14.
\end{equation}
This counting argument applies to the dependent stage lengths generated by the adaptive interaction.

\noindent\textbf{From incomplete discovery to a gradient lower bound.}
On the coupling event in \Cref{lem:sequential-coupling}, an output satisfying
$q_i\ge d_2$ for every $i$ requires every stage to have been discovered.
Otherwise, either the current stage fails its discovery test, giving
$q_i\le d_0<d_2$, or a later stage obeys the future-flat bound
\eqref{eq:future-good}, again giving $q_i<d_2$.
Combining~\eqref{eq:all-stages-prob} with the coupling error, for every
$N<Tn_0/2$ the final output therefore satisfies $\min_iq_i<d_2$ with
probability at least
\[
 1-\frac14-\frac1{12}=\frac23.
\]
The stationarity barrier in \Cref{lem:stationarity-barrier} yields
\[
 \E_{U,\xi}\norm{\nabla F_U(\widehat v)}
 \ge \frac23\,4\eps>\eps.
\]

\noindent\textbf{Randomization and the final rate.}
The same distribution and dimension work for every deterministic policy.
Conditioning a randomized algorithm on its independent seed and then averaging
therefore preserves this strict lower bound.  Consequently, some fixed frame
$U$ satisfies $\E_{R,\xi}\|\nabla F_U(\widehat v)\|>\eps$, as detailed in
\Cref{app:yao}.  Finally, the nontrivial-regime condition and the choice of
$c_T$ imply
$T=\Theta(\Delta\Lb/\eps^2)$, while~\eqref{eq:n0} gives
$n_0=\Theta(1+\sigma^2/\eps^2)$.  Hence
\[
 Tn_0
 =\Theta\!\left(
   \frac{\Delta\Lb}{\eps^2}
   \left(1+\frac{\sigma^2}{\eps^2}\right)
 \right),
\]
which is equivalent up to universal constants to~\eqref{eq:lower-rate}.
For every horizon $N$, the hard ambient dimension $d=(T+1)D$ is finite.
\end{proof}

Combining \Cref{thm:lower} with \Cref{lem:sgd} proves
\Cref{thm:main}.

\section{Discussion}\label{sec:discussion}

The sharp single-query rate separates regularity of the sample losses from
access to a shared sample. Mean-squared smoothness guarantees an accurate
common-sample gradient difference, but the single-query oracle does not
provide that difference. The lower bound shows that this restriction has a
minimax cost even when every sample gradient is globally Lipschitz.
Together with the classical upper bound, it establishes the optimality of
SGD in both the mean-squared-smooth class and its samplewise-smooth
subclass, under the same deterministic query cap and expected-stationarity
criterion.

Two features of the construction explain the stochastic factor. Additive
noise preserves sample smoothness without increasing the variance as the
number of stages grows. The sequential potential then requires the
algorithm to learn a new direction at every stage. Because the oracle
returns gradients of scalar losses, it also exposes exact information in
other variable blocks. The proof accounts for this information through
adaptive scalar measurements rather than treating the noisy vector as
the entire transcript. The resulting information cost can therefore be
summed across stages for unrestricted adaptive queries.

The characterization is dimension-uniform and uses only a variance bound
on the noise. Its finite-dimensional hard instances may grow with the
query horizon, and their Gaussian perturbations have unbounded support.
Dependence on a prescribed dimension or confidence level, and guarantees
formulated using expected stopping budgets, are outside the stated
characterization. Another question is whether the same rate can be
realized by samplewise-smooth losses with uniformly bounded gradient
noise. The present construction leaves this joint restriction unresolved.

\appendix
\section{Elementary Probability Facts}\label[appendix]{app:probability}

\begin{lemma}[Spherical cap]\label[lemma]{lem:spherical-cap}
Let $u$ be uniform on the unit sphere of an $m$-dimensional Euclidean space.
For every fixed vector $v$ with $\|v\|\le1$ and every $t\in(0,1)$,
\[
 \Pp(|\ip{u}{v}|\ge t)
 \le 2\exp\!\left(-\frac{(m-1)t^2}{2}\right).
\]
The same bound holds conditionally whenever $u$ is conditionally uniform on a
specified $m$-dimensional residual subspace and $v$ is measurable before $u$'s
remaining orientation is exposed.
\end{lemma}
\begin{proof}
Project $v$ onto the sphere's ambient subspace. Its norm is at most one,
so rotational invariance reduces the claim to $v=e_1$. Write
$u=g/\|g\|$ with $g\sim N(0,I_m)$, and set
$S^2=\sum_{i=2}^m g_i^2$. The variables $g_1$ and $S$ are independent.
The Gaussian tail bound and the chi-square Laplace transform give
\begin{align*}
 \Pp(|u_1|\ge t)
 &=\Pp\!\left(|g_1|\ge\frac{tS}{\sqrt{1-t^2}}\right)\\
 &\le 2\E\exp\!\left(-\frac{t^2S^2}{2(1-t^2)}\right)\\
 &=2(1-t^2)^{(m-1)/2}
 \le2\exp\!\left(-\frac{(m-1)t^2}{2}\right).
\end{align*}
The conditional statement is the same calculation in the specified
residual subspace.
\end{proof}

\begin{lemma}[Chi-square concentration]\label[lemma]{lem:chi-square}
If $X\sim\chi_k^2$, then for every $x\ge4k$,
\[
 \Pp(X\ge x)\le e^{-x/8}.
\]
If $Y\sim\chi_m^2$, then
\[
 \Pp(Y\le m/2)\le e^{-m/16}.
\]
\end{lemma}
\begin{proof}
Both inequalities follow from the chi-square moment-generating function and Chernoff's method. For the upper tail, $t=1/4$ gives $\Pp(X\ge x)\le 2^{k/2}e^{-x/4}\le e^{-x/8}$ when $x\ge4k$. For the lower tail, $\E e^{-Y/2}=2^{-m/2}$ gives $\Pp(Y\le m/2)\le e^{m/4}2^{-m/2}\le e^{-m/16}$.
\end{proof}

\begin{lemma}[Binary data processing bound]\label[lemma]{lem:binary-kl}
For probability measures $P,Q$ and an event $A$, writing
$p=P(A)$ and $q=Q(A)$,
\[
 \KL(P\|Q)\ge \operatorname{kl}(p\|q)
 \ge p\log(1/q)-\log2.
\]
\end{lemma}
\begin{proof}
The first inequality is data processing under the map $\omega\mapsto
\mathbf1_A(\omega)$.  For the second, expand binary relative entropy and use
$p\log p+(1-p)\log(1-p)\ge-\log2$
and $(1-p)\log(1-q)\le0$.
\end{proof}

\section{The Fixed-Instance Randomized Conclusion}\label[appendix]{app:yao}
The lower-bound proof first fixes a deterministic query rule and averages
over the Haar frame. The hard dimension and all construction parameters
depend only on the public parameters and the proposed query cap. They are
therefore the same for every realization of an algorithm's random seed.

Let a randomized algorithm have seed $R$, independent of the frame and the
oracle noise. Conditioning on $R=r$ gives a deterministic rule to which
the same distributional lower bound applies. Tonelli's theorem and the
uniform strict margin in~\Cref{thm:lower} imply
\[
 \E_U\E_{R,\xi}\|\nabla F_U(\widehat v)\|>\eps.
\]
Consequently, at least one fixed frame $U$ satisfies
$\E_{R,\xi}\|\nabla F_U(\widehat v)\|>\eps$.  Thus the distributional argument
yields the worst-case randomized lower bound in~\Cref{def:complexity}.
This averaging step fixes the entire instance before the algorithm's
internal randomness and oracle samples are drawn.

\begin{thebibliography}{99}
\bibitem{ghadimi}
S. Ghadimi and G. Lan.
\newblock Stochastic first- and zeroth-order methods for nonconvex stochastic programming.
\newblock \emph{SIAM Journal on Optimization}, 23(4):2341--2368, 2013.
\newblock \url{https://doi.org/10.1137/120880811}.

\bibitem{spider}
C. Fang, C. J. Li, Z. Lin, and T. Zhang.
\newblock SPIDER: Near-optimal non-convex optimization via stochastic path-integrated differential estimator.
\newblock \emph{Advances in Neural Information Processing Systems}, 31, 2018.
\newblock \url{https://arxiv.org/abs/1807.01695}.

\bibitem{arjevani}
Y. Arjevani, Y. Carmon, J. C. Duchi, D. J. Foster, N. Srebro, and B. E. Woodworth.
\newblock Lower bounds for non-convex stochastic optimization.
\newblock \emph{Mathematical Programming}, 199:165--214, 2023.
\newblock \url{https://doi.org/10.1007/s10107-022-01822-7}.

\bibitem{snvrg}
D. Zhou, P. Xu, and Q. Gu.
\newblock Stochastic nested variance reduction for nonconvex optimization.
\newblock \emph{Advances in Neural Information Processing Systems}, 31, 2018.
\newblock \url{https://arxiv.org/abs/1806.07811}.

\bibitem{storm}
A. Cutkosky and F. Orabona.
\newblock Momentum-based variance reduction in non-convex SGD.
\newblock \emph{Advances in Neural Information Processing Systems}, 32, 2019.
\newblock \url{https://proceedings.neurips.cc/paper/2019/hash/b8002139cdde66b87638f7f91d169d96-Abstract.html}.

\bibitem{page}
Z. Li, H. Bao, X. Zhang, and P. Richt\'arik.
\newblock PAGE: A simple and optimal probabilistic gradient estimator for nonconvex optimization.
\newblock \emph{Proceedings of the 38th International Conference on Machine Learning}, PMLR 139:6286--6295, 2021.
\newblock \url{https://proceedings.mlr.press/v139/li21a.html}.

\bibitem{fradin2026}
A. Fradin, A. Sadiev, L. Condat, and P. Richt\'arik.
\newblock Tight lower bounds and optimal algorithms for stochastic nonconvex optimization with heavy-tailed noise.
\newblock \emph{Proceedings of the 29th International Conference on Artificial Intelligence and Statistics}, PMLR 300:2548--2556, 2026.
\newblock \url{https://proceedings.mlr.press/v300/fradin26a.html}.
\newblock Extended version: \url{https://arxiv.org/abs/2512.18713v2}.

\bibitem{carmon}
Y. Carmon, J. C. Duchi, O. Hinder, and A. Sidford.
\newblock Lower bounds for finding stationary points I.
\newblock \emph{Mathematical Programming}, 184:71--120, 2020.
\newblock \url{https://arxiv.org/abs/1710.11606}.

\bibitem{peng2026}
Y. Peng, Z. Tang, and W. Jia.
\newblock Dense weak hiding: Closing complexity gaps in nonconvex and PL finite-sum optimization under individual smoothness.
\newblock arXiv preprint arXiv:2609.00045, version 3, 2026.
\newblock \url{https://arxiv.org/abs/2609.00045v3}.

\bibitem{jin2026}
J. Jin.
\newblock A tight lower bound for smooth nonconvex stochastic optimization with bounded gradient noise.
\newblock arXiv preprint arXiv:2608.09004, 2026.
\newblock \url{https://arxiv.org/abs/2608.09004}.

\bibitem{secondorder}
Y. Arjevani, Y. Carmon, J. C. Duchi, D. J. Foster, A. Sekhari, and K. Sridharan.
\newblock Second-order information in non-convex stochastic optimization: Power and limitations.
\newblock \emph{Proceedings of the 33rd Conference on Learning Theory}, PMLR 125:242--299, 2020.
\newblock \url{https://proceedings.mlr.press/v125/arjevani20a.html}.
\end{thebibliography}

\end{document}
